% GENERATED FILE. Edit canonical lessons, then run npm run generate:books.
% canonical-source-hash: fnv1a64:d4264ddacefa3f87
% canonical-lessons: TA-C246-kontattam, TA-C246-pattacu, TA-C246-kalaccaram, TA-C246-palakkam, TA-C246-parampariyam

\chapter{Celebration, Firecrackers, Culture, Habit, Tradition}
\label{ch:ta-celebration-firecrackers-culture-habit-tradition}
\hlchaptermodality{246}

\begin{chapteropening}
\textbf{By the end of this chapter:} \emph{I can say celebration, firecrackers, culture, a habit and tradition.}

Complete the last lesson of chapter 246: I can say celebration, firecrackers, culture, a habit and tradition.
\end{chapteropening}

\section[\texorpdfstring{\ta{கொண்டாட்டம்} (koṇṭāṭṭam)}{கொண்டாட்டம் (koṇṭāṭṭam)}]{\ta{கொண்டாட்டம்} (koṇṭāṭṭam) — celebration, festivities}

\label{lesson:TA-C246-kontattam}



\begin{quote}
Before the new one: say the Tamil for a man, then the Tamil for a woman.
\end{quote}

\subsection*{You'll want to know: \ta{கொண்டாட்டம்}}
\textbf{\ta{கொண்டாட்டம்}} — \emph{koṇṭāṭṭam} — \textquotedblleft{}celebration, festivities\textquotedblright{}.

\begin{cousinweb}
\textbf{\ta{கொண்டாட்டம்}} (\emph{koṇṭāṭṭam}) is the noun of \emph{koṇṭāṭu}, \textquotedblleft{}to celebrate\textquotedblright{}, usually analysed as native \emph{koṇṭu}, \textquotedblleft{}having taken\textquotedblright{}, from \emph{koḷ} (DEDR 2151), and \emph{āṭu}, \textquotedblleft{}to dance, to play\textquotedblright{} (DEDR 347).
\end{cousinweb}

\begin{grammarlens}[title={What you've built}]
One more word for this chapter.
\end{grammarlens}

\subsection*{Your turn}
\begin{itemize}
\raggedright
  \item \emph{Say it:} \emph{koṇṭāṭṭam}
  \item \emph{Say it:} \emph{koṇṭāṭṭam}, once more
  \item \emph{Say it:} say \emph{koṇṭāṭṭam}
  \item \emph{Recall:} say the Tamil for a man, then the Tamil for a woman, then say \emph{koṇṭāṭṭam} again
\end{itemize}

\begin{tcolorbox}[breakable,colback=teal!4,colframe=teal!35!black,title={Before you move on}]
What does \ta{கொண்டாட்டம்} mean? (\textquotedblleft{}Celebration, festivities\textquotedblright{}.) Say it once more.
\end{tcolorbox}
\section[\texorpdfstring{\ta{பட்டாசு} (paṭṭācu)}{பட்டாசு (paṭṭācu)}]{\ta{பட்டாசு} (paṭṭācu) — firecrackers}

\label{lesson:TA-C246-pattacu}



\begin{quote}
Before the new one: say the Tamil for a woman, then the Tamil for celebration.
\end{quote}

\subsection*{You'll want to know: \ta{பட்டாசு}}
\textbf{\ta{பட்டாசு}} — \emph{paṭṭācu} — \textquotedblleft{}firecrackers\textquotedblright{}.

\begin{cousinweb}
The origin of \textbf{\ta{பட்டாசு}} (\emph{paṭṭācu}) is uncertain. DEDR has no entry for it; it is often compared with Hindustani \emph{paṭākhā}, \textquotedblleft{}firecracker\textquotedblright{}.
\end{cousinweb}

\begin{grammarlens}[title={What you've built}]
One more word for this chapter.
\end{grammarlens}

\subsection*{Your turn}
\begin{itemize}
\raggedright
  \item \emph{Say it:} \emph{paṭṭācu}
  \item \emph{Say it:} \emph{paṭṭācu}, once more
  \item \emph{Say it:} say \emph{paṭṭācu}
  \item \emph{Recall:} say the Tamil for a woman, then the Tamil for celebration, then say \emph{paṭṭācu} again
\end{itemize}

\begin{tcolorbox}[breakable,colback=teal!4,colframe=teal!35!black,title={Before you move on}]
What does \ta{பட்டாசு} mean? (\textquotedblleft{}Firecrackers\textquotedblright{}.) Say it once more.
\end{tcolorbox}
\section[\texorpdfstring{\ta{கலாச்சாரம்} (kalāccāram)}{கலாச்சாரம் (kalāccāram)}]{\ta{கலாச்சாரம்} (kalāccāram) — culture}

\label{lesson:TA-C246-kalaccaram}



\begin{quote}
Before the new one: say the Tamil for celebration, then the Tamil for firecrackers.
\end{quote}

\subsection*{You'll want to know: \ta{கலாச்சாரம்}}
\textbf{\ta{கலாச்சாரம்}} — \emph{kalāccāram} — \textquotedblleft{}culture\textquotedblright{}.

\begin{cousinweb}
\textbf{\ta{கலாச்சாரம்}} (\emph{kalāccāram}) is a modern coinage from Sanskrit \emph{kalā}, \textquotedblleft{}art\textquotedblright{} (CDIAL 2927), and \emph{ācāra}, \textquotedblleft{}conduct, custom\textquotedblright{} (CDIAL 1071). Which way \emph{kalā} travelled is debated: CDIAL takes it from Dravidian, while DEDR files the native Tamil word for \textquotedblleft{}art\textquotedblright{} under the verb \textquotedblleft{}to learn\textquotedblright{} (DEDR 1297) and only compares \emph{kalā}.
\end{cousinweb}

\begin{grammarlens}[title={What you've built}]
One more word for this chapter.
\end{grammarlens}

\subsection*{Your turn}
\begin{itemize}
\raggedright
  \item \emph{Say it:} \emph{kalāccāram}
  \item \emph{Say it:} \emph{kalāccāram}, once more
  \item \emph{Say it:} say \emph{kalāccāram}
  \item \emph{Recall:} say the Tamil for celebration, then the Tamil for firecrackers, then say \emph{kalāccāram} again
\end{itemize}

\begin{tcolorbox}[breakable,colback=teal!4,colframe=teal!35!black,title={Before you move on}]
What does \ta{கலாச்சாரம்} mean? (\textquotedblleft{}Culture\textquotedblright{}.) Say it once more.
\end{tcolorbox}
\section[\texorpdfstring{\ta{பழக்கம்} (paḻakkam)}{பழக்கம் (paḻakkam)}]{\ta{பழக்கம்} (paḻakkam) — a habit, a custom}

\label{lesson:TA-C246-palakkam}



\begin{quote}
Before the new one: say the Tamil for firecrackers, then the Tamil for culture.
\end{quote}

\subsection*{You'll want to know: \ta{பழக்கம்}}
\textbf{\ta{பழக்கம்}} — \emph{paḻakkam} — \textquotedblleft{}a habit, a custom\textquotedblright{}.

\begin{cousinweb}
\textbf{\ta{பழக்கம்}} (\emph{paḻakkam}) is native Dravidian. DEDR 4000 lists Malayalam \emph{paḻakuka}, \textquotedblleft{}to be accustomed\textquotedblright{}. It is the noun of \emph{paḻaku}, \textquotedblleft{}to become used to\textquotedblright{}.
\end{cousinweb}

\begin{grammarlens}[title={What you've built}]
One more word for this chapter.
\end{grammarlens}

\subsection*{Your turn}
\begin{itemize}
\raggedright
  \item \emph{Say it:} \emph{paḻakkam}
  \item \emph{Say it:} \emph{paḻakkam}, once more
  \item \emph{Say it:} say \emph{paḻakkam}
  \item \emph{Recall:} say the Tamil for firecrackers, then the Tamil for culture, then say \emph{paḻakkam} again
\end{itemize}

\begin{tcolorbox}[breakable,colback=teal!4,colframe=teal!35!black,title={Before you move on}]
What does \ta{பழக்கம்} mean? (\textquotedblleft{}A habit, a custom\textquotedblright{}.) Say it once more.
\end{tcolorbox}
\section[\texorpdfstring{\ta{பாரம்பரியம்} (pārampariyam)}{பாரம்பரியம் (pārampariyam)}]{\ta{பாரம்பரியம்} (pārampariyam) — tradition, heritage}

\label{lesson:TA-C246-parampariyam}



\begin{quote}
Before the new one: say the Tamil for culture, then the Tamil for a habit.
\end{quote}

\subsection*{You'll want to know: \ta{பாரம்பரியம்}}
\textbf{\ta{பாரம்பரியம்}} — \emph{pārampariyam} — \textquotedblleft{}tradition, heritage\textquotedblright{}.

\begin{cousinweb}
\textbf{\ta{பாரம்பரியம்}} (\emph{pārampariyam}) is a well-established loan from Sanskrit \emph{pāramparya}, \textquotedblleft{}succession, handing down\textquotedblright{}, built on \emph{paramparā}, \textquotedblleft{}an unbroken series\textquotedblright{}.
\end{cousinweb}

\begin{grammarlens}[title={What you've built}]
One more word for this chapter.
\end{grammarlens}

\subsection*{Your turn}
\begin{itemize}
\raggedright
  \item \emph{Say it:} \emph{pārampariyam}
  \item \emph{Say it:} \emph{pārampariyam}, once more
  \item \emph{Say it:} say \emph{pārampariyam}
  \item \emph{Recall:} say the Tamil for culture, then the Tamil for a habit, then say \emph{pārampariyam} again
\end{itemize}

\begin{tcolorbox}[breakable,colback=teal!4,colframe=teal!35!black,title={Before you move on}]
What does \ta{பாரம்பரியம்} mean? (\textquotedblleft{}Tradition, heritage\textquotedblright{}.) Say it once more.
\end{tcolorbox}
